\documentclass[a4paper]{article}
\usepackage[pdftex]{graphicx}
\usepackage[utf8]{inputenc}
\usepackage{enumerate}
\usepackage{amssymb}
\usepackage{icomma}
\usepackage{siunitx}
\sisetup{locale=DE} 
\usepackage{tikz}
\usepackage{href-ul}
\hypersetup{
	colorlinks=true,
	linkcolor=blue,
	urlcolor=blue}
\usepackage{eurosym}
\usepackage{geometry}
\geometry{a4paper, top=15mm, left=15mm, right=15mm, bottom=15mm,
	headsep=10mm, footskip=12mm}

\begin{document}
\begin{tikzpicture}(10,3)
	\draw[ultra thick](2,0) --(10,0) -- (10,1.3) --(2,1.3) -- (2,0);
	\draw[fill=black](2,0)-- (0,.6) -- (2,.6) -- (2,0);
	\node at (6,.6) {\large Lösungsblatt von \href{https://www.okuyakl.de}{www.okuyakl.de}};
\end{tikzpicture}
\vspace{0.5 cm}

\noindent{\bf Aufgabe 1. a)}\\
$$25\% \qquad 80\% \qquad 20\% \qquad 28\% \qquad 75\% $$

\noindent{\bf Aufgabe 1. b)}\\
$$\SI{21}{\meter} \qquad \SI{51}{\kilogram} \qquad \textnormal{\EUR{60}} \qquad \SI{180}{\milli\gram} $$

\noindent{\bf Aufgabe 2.}\\
Wenn \EUR{50} = 100\% sind, dann sind \EUR{21} = 42\%. Die Ware wurde also um 58\% reduziert, nicht um 65\%.

\noindent{\bf Aufgabe 3.}\\
Die Prozentformel für den Prozentsatz lautet (G = Grund\-wert; P = Prozentwert):
$$ p\% = {P \over G} \cdot 100\% $$

\noindent{\bf Aufgabe 4. a)}\\
Der Endbetrag ist
$$\textnormal{\EUR{1280}} \cdot 1,19 = \textnormal{\EUR{1523,20}}$$

\noindent{\bf Aufgabe 4. b)}\\
Der Endbetrag ist
$$\textnormal{\EUR{1523,20}} \cdot 0,98 = \textnormal{\EUR{1492,74}} $$

\noindent{\bf Aufgabe 5.}\\
430 ist mehr als die Hälfte von 855, damit ist der Antrag angenommen.

\noindent{\bf Aufgabe 6.}\\
Die Mäntel kosten:
$$\textnormal{\EUR{119}} \cdot 0,78 = \textnormal{\EUR{92,82}} $$
Die Anzüge kosten:
$$\textnormal{\EUR{169}} \cdot 0,70 = \textnormal{\EUR{118,30}} $$

\noindent{\bf Aufgabe 7.}\\
$$Z_t = {K \cdot p \cdot t \over 100 \cdot 360}$$

\noindent{\bf Aufgabe 8.}\\
Die Jahreszinsen müssen betragen:
$$Z = \textnormal{\EUR{600}} \cdot 12 = \textnormal{\EUR{7200}}$$
Das hierfür verzinste Kapital muss sein:
$$ K = {100\% \over 6 \%} \cdot \textnormal{\EUR{7200}} = \textnormal{\EUR{120000}}$$

\noindent{\bf Aufgabe 9.}\\
Die Jahreszinsen sind:
$$Z = \textnormal{\EUR{375}} \cdot 12 = \textnormal{\EUR{4500}}$$
Der Zinssatz ist:
$$ p = {\textnormal{\EUR{4500}} \over \textnormal{\EUR{90000}}}\cdot 100\% = 5\%$$

\noindent{\bf Aufgabe 10. a)}\\
Zinsen aus (1): \EUR{2700} \qquad Zinsen aus (2): \EUR{3920} \qquad Zinsen aus (3): \EUR{4880}\\
Zinsen insgesamt:  \EUR{11500} \\

\noindent{\bf Aufgabe 10. b)}\\
Die Tilgung ist:
$$(50000+70000+80000)\cdot 0,04 = \textnormal{\EUR{8000}}$$
Die Gesamtbelastung ist:
$$11500+8000= \textnormal{\EUR{19500}}$$

\noindent{\bf Aufgabe 11.}\\
Das Wasservolumen ist:
$$\SI{50}{\liter} \cdot 0,91 = \SI{45,5}{\liter}$$

\noindent{\bf Aufgabe 12.}\\
Leitung A füllt ${1\over 20}$ des Beckens pro Stunde, Leitung B ${1 \over 12}$.
Zusammen füllen sie pro Stunde:
$${1\over 20}+ {1 \over 12} ={2 \over 15}$$
Der Kehrwert hiervon ist die Dauer T, bis das ganze Becken voll ist:
$$T= {15 \over 2}=\SI{7,5}{\hour}$$

\begin{center}
	\includegraphics[width=7 cm]{../../viecher/endcomic.pdf}
	
	Hier geht es zurück zum \href{https://www.okuyakl.de/math/m7proL055/aa055.pdf}{Aufgabenblatt}
\end{center}


\end{document}

